%! Rates of Change — Unit 1, Lesson 1.3 — Guided Notes (Answer Key)
\documentclass[10pt]{article}
\usepackage{precalculus-article}
\usepackage{precalculus-key}   % pulls in -boxes; do NOT also load -boxes
\newcommand{\TallMath}[1]{$\displaystyle #1\rule[-1.4em]{0pt}{3.2em}$}

\begin{document}

\pageheader{Unit 1, Lesson 1.3}{Guided Notes --- Answer Key}

\vspace{0.12in}

\begin{objectivebox}
By the end of this lesson, I will be able to\ldots
\begin{enumerate}[itemsep=2pt]
  \item Compute the \ans{average rate of change} of a function
        over an interval from a formula, a table, or a graph.
  \item Interpret the \ans{sign} and \ans{units} of a rate of change in
        context.
  \item Use the pattern of AROCs to tell \ans{linear} functions from
        \ans{quadratic} ones, and to name concavity.
\end{enumerate}
\end{objectivebox}

\vspace{0.1in}

\boxguard
\begin{vocabbox}
Fill in each term as we define it together.
\par\vspace{2pt}
\noindent\textbf{\textcolor{plum}{Rate of change:}}\\[1pt]
\ansline{How fast the output changes with respect to the input --- output units per input unit.}\\[3pt]
\noindent\textbf{\textcolor{plum}{Average rate of change (AROC):}}\\[1pt]
\ansline{Over $[a,b]$: the ratio $(f(b)-f(a))/(b-a)$ --- the constant rate producing the same net change.}\\[3pt]
\noindent\textbf{\textcolor{plum}{Secant line:}}\\[1pt]
\ansline{The line through two points $(a, f(a))$ and $(b, f(b))$ on a graph; its slope is the AROC.}\\[3pt]
\noindent\textbf{\textcolor{plum}{Positive / negative rate of change:}}\\[1pt]
\ansline{Positive: the quantities move in the same direction. Negative: opposite directions.}\\[3pt]
\noindent\textbf{\textcolor{plum}{First differences:}}\\[1pt]
\ansline{The AROCs over consecutive equal-length intervals --- the ``steps'' in an equally spaced table.}
\end{vocabbox}

\vspace{0.1in}

\boxguard
\begin{hookbox}
A family drives 120 miles in 3 hours, and the GPS reports:
\textit{average speed 40 mph}. They were certainly not going exactly 40 the
whole time --- so what does the 40 \textit{promise}?

\ansline{One steady speed of 40 mph would have covered the same 120 miles in the same 3 hours.}
\ansline{The average is the one constant rate that produces the same net change.}
\end{hookbox}

\vspace{0.1in}

\boxguard[34]
\begin{notesbox}{1. Average Rate of Change --- Slope, Upgraded}
The slope formula measures steepness between two points. The
\textbf{average rate of change} is the same fraction, written with function
notation --- and it works on \textit{any} function, not just lines.

\begin{center}
$\text{AROC over } [a, b] \;=\;
\dfrac{\text{change in \ans{output}}}{\text{change in \ans{input}}}
\;=\; \dfrac{f(b) - f(a)}{b - a}$
\end{center}

\textbf{Example 1.} Let $f(x) = x^2$. Find the AROC over $[1, 4]$.

\begin{work}
  f(1) &= 1 \\
  f(4) &= 16 \\
  \text{AROC} &= \frac{16 - 1}{4 - 1} \\
              &= \frac{15}{3} \\
              &= 5
\end{work}

\begin{itemize}[itemsep=4pt]
  \item The line through $(1, 1)$ and $(4, 16)$ is called a
        \ans{secant} line, and the AROC is its \ans{slope}.
  \item Meaning: 5 is the one \ans{constant} rate that would produce the same
        net change ($+15$) over the same interval --- the GPS promise.
\end{itemize}
\end{notesbox}

\vspace{0.1in}

\boxguard
\begin{notesbox}{2. AROC from a Table and a Graph}
\textbf{From a table.} $T(t)$ is the temperature in $^\circ$F $t$ hours after
sunrise.

\smallskip
\renewcommand{\arraystretch}{1.4}
\begin{tabular}{|c|c|c|c|c|}
\hline
$t$ (hours) & 0 & 2 & 4 & 6 \\
\hline
$T(t)$ ($^\circ$F) & 40 & 52 & 58 & 54 \\
\hline
\end{tabular}
\smallskip

\begin{itemize}[itemsep=4pt]
  \item AROC over $[0, 2]$: \ans{6} $^\circ$F per hour \quad
        AROC over $[2, 4]$: \ans{3} $^\circ$F per hour
  \item AROC over $[4, 6]$: \ans{$-2$} $^\circ$F per hour \quad
        Fastest warming stretch: \ans{$[0,2]$}
\end{itemize}

\textbf{From a graph.} The graph of $f$ is below; the dashed line is the
secant through $(0, 2)$ and $(4, 6)$.

\begin{center}
\begin{tikzpicture}[x=0.6cm, y=0.5cm]
  \draw[linegray, very thin, step=1] (0,0) grid (8,7);
  \draw[->, thick] (0,0) -- (8.6,0) node[below right] {\small $x$};
  \draw[->, thick] (0,0) -- (0,7.6) node[above] {\small $f(x)$};
  \foreach \x in {1,2,3,4,5,6,7,8} \node[below] at (\x,0) {\small \x};
  \foreach \y in {1,2,3,4,5,6,7} \node[left] at (0,\y) {\small \y};
  \draw[very thick, plum] plot [smooth] coordinates
    {(0,2) (2,5) (4,6) (6,5) (8,2)};
  \draw[dashed, thick, goldacc] (0,2) -- (4,6);
  \fill[plum] (0,2) circle (2.5pt);
  \fill[plum] (4,6) circle (2.5pt);
  \fill[plum] (8,2) circle (2.5pt);
\end{tikzpicture}
\end{center}

\begin{itemize}[itemsep=4pt]
  \item AROC over $[0, 4]$: read both points, then
        $\dfrac{6 - 2}{4 - 0} = $ \ans{1}
  \item AROC over $[4, 8]$: \quad $\dfrac{2 - 6}{8 - 4} = $ \ans{$-1$}
\end{itemize}

\textbf{Watch the order:} output difference upstairs, input difference
downstairs --- \textit{subtracted in the same order on both floors}.
\end{notesbox}

\vspace{0.1in}

\boxguard
\begin{notesbox}{3. What the Sign Says}
\begin{itemize}[itemsep=4pt]
  \item A \textbf{positive} AROC: as one quantity increases, the other
        \ans{increases} --- the quantities move in the \ans{same} direction.
  \item A \textbf{negative} AROC: as one quantity increases, the other
        \ans{decreases} --- opposite directions.
  \item Every rate carries \textbf{units}: output units per input unit. The
        AROC over $[4,6]$ above is $-2$ \ans{degrees ($^\circ$F) per hour} , which
        says the temperature was \ans{falling} by 2 degrees each hour.
\end{itemize}
\end{notesbox}

\vspace{0.1in}

\boxguard[30]
\begin{notesbox}{4. The Pattern of the AROCs --- Linear vs.\ Quadratic}
Compute the AROC over each unit interval and watch the \textit{pattern}.

\medskip
\textbf{Linear:} $f(x) = 2x + 1$

\smallskip
\renewcommand{\arraystretch}{1.4}
\begin{tabular}{|c|c|c|c|c|c|}
\hline
$x$ & 0 & 1 & 2 & 3 & 4 \\
\hline
$f(x)$ & 1 & 3 & 5 & 7 & 9 \\
\hline
\end{tabular}
\smallskip

AROCs over $[0,1], [1,2], [2,3], [3,4]$: \quad
\ans{2} , \ans{2} , \ans{2} , \ans{2}

For a \textbf{linear} function, the AROC is \ans{constant} over every
interval.

\medskip
\textbf{Quadratic:} $g(x) = x^2$

\smallskip
\begin{tabular}{|c|c|c|c|c|c|}
\hline
$x$ & 0 & 1 & 2 & 3 & 4 \\
\hline
$g(x)$ & 0 & 1 & 4 & 9 & 16 \\
\hline
\end{tabular}
\smallskip

AROCs over the same intervals: \quad
\ans{1} , \ans{3} , \ans{5} , \ans{7}
\qquad each step changes by \ans{2}

For a \textbf{quadratic} function, the AROCs are not constant --- they change
at a \ans{constant} rate.

\medskip
\textbf{The concavity link (from 1.2):} when the AROCs are
\textit{increasing}, the graph is concave \ans{up}; when the AROCs are
\textit{decreasing}, the graph is concave \ans{down}.
\end{notesbox}

\vspace{0.1in}

\boxguard
\begin{practicebox}
\begin{enumerate}[itemsep=8pt]
  \item Let $h(x) = x^2 - 3x$. Find the AROC over $[2, 5]$.

        \begin{work}
          h(2) &= (2)^2 - 3(2) \\
               &= -2 \\
          h(5) &= (5)^2 - 3(5) \\
               &= 10 \\
          \text{AROC} &= \frac{10 - (-2)}{5 - 2} \\
                      &= 4
        \end{work}

        AROC $= $ \ans{4}

  \item A function's AROCs over consecutive equal-length intervals are
        $10,\ 7,\ 4,\ 1$.

        \begin{enumerate}[label=(\alph*), itemsep=4pt]
          \item Each AROC changes by \ans{$-3$} --- a constant step, so the
                function is \ans{quadratic}.
          \item The AROCs are decreasing, so the graph is concave
                \ans{down}.
        \end{enumerate}
\end{enumerate}
\end{practicebox}

\end{document}
